% Part: methods
% Chapter: induction
% Section: relations

\documentclass[../../../include/open-logic-section]{subfiles}

\begin{document}

\olfileid{mth}{ind}{rel}

\olsection{تعلق تے فنکشن}

جدوں اسیں شےواں دا کوئی سیٹ (جیویں قدرتی عدد یا ٹھیک ٹرم) استقرائی طریقے نال تعریف کر لئیے، تاں ایہناں شےواں \emph{اُتے تعلقاں} دی تعریف وی استقرا نال دے سکدے آں۔ مثلاً ایہ خیال ویکھو۔ اک ٹھیک ٹرم~$t_1$، اک ٹھیک ٹرم~$t_2$ دا ذیلی ٹرم اے جے اوہ اوہدے اک حصے وجوں آؤندا اے۔ ایس لئی اک نشان ورتیے، $t_1 \sqsubseteq t_2$۔ بے شک ہر ٹھیک ٹرم اپنا ذیلی ٹرم اے، $t \sqsubseteq t$۔ ایس تعلق دی استقرائی تعریف ایویں دے سکدے آں:

\begin{defn}
  ایہ تعلق کہ ٹھیک ٹرم~$t_1$،~$t_2$ دا ذیلی ٹرم اے، $t_1
  \sqsubseteq t_2$،~$t_2$ اُتے استقرا نال ایویں تعریف ہُندا اے:
  \begin{enumerate}
  \item جے $t_2$ اکھر اے، تاں $t_1 \sqsubseteq t_2$ اُدوں تے صرف اُدوں جدوں $t_1 = t_2$۔
  \item جے $t_2$، $[s_1 \circ s_2]$ اے، تاں $t_1 \sqsubseteq t_2$ اُدوں تے صرف اُدوں جدوں $t_1 =
    t_2$، یا $t_1 \sqsubseteq s_1$، یا $t_1 \sqsubseteq s_2$۔
  \end{enumerate}
\end{defn}

مثلاً ایہ تعریف دسدی اے کہ $\mathrm{a}
\sqsubseteq [\mathrm{b} \circ \mathrm{a}]$۔ کیوں جے (۲) دسدا اے کہ
$\mathrm{a} \sqsubseteq [\mathrm{b} \circ \mathrm{a}]$ اُدوں تے صرف اُدوں جدوں $\mathrm{a}
= [\mathrm{b} \circ \mathrm{a}]$، یا $\mathrm{a} \sqsubseteq b$، یا
$\mathrm{a} \sqsubseteq \mathrm{a}$۔ پہلیاں دو گلاں غلط نیں۔ صاف اے کہ $\mathrm{a}$، $[\mathrm{b} \circ
  \mathrm{a}]$ دے اکّو نہیں، تے~(۱) مطابق $\mathrm{a} \sqsubseteq \mathrm{b}$ اُدوں تے صرف اُدوں جدوں $\mathrm{a} = \mathrm{b}$، جیہڑی گل وی غلط اے۔ پر~(۱) مطابق ای، $\mathrm{a} \sqsubseteq \mathrm{a}$ اُدوں تے صرف اُدوں جدوں $\mathrm{a} = \mathrm{a}$، جیہڑی گل سچ اے۔

ایہ دھیان رکھنا ضروری اے کہ تعریف دی کامیابی اک ایہو جیہی گل اُتے ٹکدی اے جیہڑی اسیں ہن تک ثابت نہیں کیتی۔ ہر ٹھیک ٹرم~$t$ یا بس اک اکھر اے، یا \emph{اکّو طریقے نال طے ہون والے} ٹھیک ٹرم $s_1$ تے $s_2$ نیں جیہناں لئی $t = [s_1 \circ s_2]$۔ «اکّو طریقے نال طے ہون والے» دا مطلب اے کہ جے $t = [s_1 \circ s_2]$، تاں ایہ \emph{نال ای} $= [r_1 \circ r_2]$ نہیں ہو سکدا، جیہدے وچ $s_1 \neq r_1$ یا $s_2 \neq r_2$ ہووے۔ جے ایہ ہو سکدا، تاں شق~(۲) دا اپنے آپ نال ٹکراؤ ہو سکدا سی۔ $t_2$ نوں $[s_1 \circ s_2]$ پڑھ کے شاید سانوں $t_1
\sqsubseteq t_2$ ملے، پر $t_2$ نوں $[r_1 \circ r_2]$ پڑھ کے شاید سانوں $t_1 \sqsubseteq t_2$ دی نفی ملے۔ ایہ ثابت کرن توں پہلاں کہ ایہ نہیں ہو سکدا، اک مثال ویکھیے جتھے ایہ ہو \emph{سکدا} اے۔

\begin{defn}
  \emph{بریکٹاں توں بناں ٹرم} استقرائی طریقے نال ایویں تعریف کرو:
  \begin{enumerate}
  \item ہر اکھر بریکٹاں توں بناں اک ٹرم اے۔
  \item جے $s_1$ تے $s_2$ بریکٹاں توں بناں ٹرم نیں، تاں $s_1 \circ s_2$ بریکٹاں توں بناں اک ٹرم اے۔
  \item ہور کوئی شے بریکٹاں توں بناں ٹرم نہیں۔
  \end{enumerate}
\end{defn}

بریکٹاں توں بناں ٹرماں دیاں مثالاں $\mathrm{a}$، $\mathrm{b} \circ
\mathrm{d}$، $\mathrm{b} \circ \mathrm{a} \circ \mathrm{b}$ نیں۔ ہُن جے بریکٹاں توں بناں ٹرماں لئی «ذیلی ٹرم» اُتے والے طریقے نال تعریف کریے، تاں دوجی شق ایویں ہووے گی:
\begin{center}
  جے $t_2 = s_1 \circ s_2$، تاں $t_1 \sqsubseteq t_2$ اُدوں تے صرف اُدوں جدوں $t_1 = t_2$، یا $t_1
  \sqsubseteq s_1$، یا $t_1 \sqsubseteq s_2$۔
\end{center}

ہُن $\mathrm{b} \circ \mathrm{a} \circ \mathrm{b}$، $s_1
\circ s_2$ دی صورت وچ اے، جتھے
\begin{align*}
  s_1 & = \mathrm{b} \text{ تے} & s_2 & = \mathrm{a} \circ \mathrm{b}.
\intertext{ایہ $r_1 \circ r_2$ دی صورت وچ وی اے، جتھے}
  r_1 & = \mathrm{b} \circ \mathrm{a} \text{ تے} & r_2 &= \mathrm{b}.
\end{align*}
ہُن کی $\mathrm{a} \circ \mathrm{b}$، $\mathrm{b} \circ
\mathrm{a} \circ \mathrm{b}$ دا ذیلی ٹرم اے؟ پہلی پڑھت لئی جواب ہاں اے تے دوجی لئی نہیں۔

ایس خاصیت نوں \emph{اکّو طرح پڑھے جانا} آکھدے نیں کہ ہور ٹھیک ٹرماں توں کسے ٹھیک ٹرم دی بناوٹ دا طریقہ اکّو اے۔ کیوں جے ایہو جیہیاں استقرائی طریقے نال تعریف ہوئیاں شےواں لئی تعلقاں دیاں استقرائی تعریفاں ضروری نیں، سانوں ثابت کرنا پئے گا کہ ایہ خاصیت پوری ہُندی اے۔

\begin{prop}
  منّ لو $t$ اک ٹھیک ٹرم اے۔ تاں یا $t$ بس اکلّا اکھر اے، یا اکّو طریقے نال طے ہون والے ٹھیک ٹرم $s_1$، $s_2$ نیں جیہناں لئی~$t =
  [s_1 \circ s_2]$۔
\end{prop}

\begin{proof}
  جے $t$ اکلّا اکھر اے، تاں شرط پوری اے۔ سو منّ لو $t$ اکلّا اکھر نہیں۔ استقرائی تعریف توں دس سکدے آں کہ فیر $t$ لازماً $[s_1 \circ s_2]$ دی صورت وچ اے، کسے ٹھیک ٹرماں $s_1$ تے~$s_2$ لئی۔ ہن بس ایہ وکھاؤنا اے کہ ایہ دوویں اکّو طریقے نال طے نیں، یعنی جے $t = [r_1 \circ r_2]$، تاں $s_1 = r_1$ تے $s_2 = r_2$۔

  سو منّ لو $t = [s_1 \circ s_2]$ تے $t = [r_1 \circ r_2]$ وی، کسے ٹھیک ٹرماں $s_1$، $s_2$، $r_1$، $r_2$ لئی۔ سانوں وکھاؤنا اے کہ $s_1 = r_1$ تے $s_2 = r_2$۔ پہلاں، $s_1$ تے $r_1$ لازماً اکّو نیں، کیوں جے نہیں تاں اک دوجے دا حقیقی مُڈھلا ٹکڑا اے۔ پر \olref[sti]{prop:initial} مطابق جے $s_1$ تے~$r_1$ دوویں ٹھیک ٹرم نیں تاں ایہ ناممکن اے۔ پر جے $s_1 = r_1$، تاں صاف اے کہ $s_2 = r_2$ وی۔
\end{proof}

اسیں فنکشن وی استقرائی طریقے نال تعریف کر سکدے آں۔ مثلاً فنکشن~$f$ دی تعریف ایویں دے سکدے آں جیہڑا کسے وی ٹھیک ٹرم نوں اوہدے اندر اک دوجے وچ پئے $[\dots]$ دی سبھ توں ودھ ڈونگھائی ول بھیجدا اے:

\begin{defn}
  \ollabel{defn:depth} ٹھیک ٹرم دی \emph{ڈونگھائی}، $f(t)$، استقرائی طریقے نال ایویں تعریف ہُندی اے:
  \[
  f(t) = \begin{cases}
    0 & \text{ جے $t$ اک اکھر اے}\\
    \max(f(s_1), f(s_2)) + 1 & \text{ جے $t = [s_1 \circ s_2]$۔}
  \end{cases}
  \]
\end{defn}

مثلاً
\begin{align*}
  f([\mathrm{a} \circ \mathrm{b}]) & =
    \max(f(\mathrm{a}),f(\mathrm{b})) + 1 = \\
  &= \max(0, 0) + 1 = 1, \text{ تے}\\
  f([[\mathrm{a} \circ \mathrm{b}] \circ \mathrm{c}]) & =
    \max(f([\mathrm{a} \circ \mathrm{b}]), f(\mathrm{c})) + 1 = \\
  & = \max(1,0) + 1 = 2.
\end{align*}

بے شک ایتھے اسیں منّدے آں کہ $s_1$ تے $s_2$ ٹھیک ٹرم نیں، تے ایہ گل ورتدے آں کہ ہر ٹھیک ٹرم یا اکھر اے یا $[s_1 \circ s_2]$ دی صورت وچ اے۔ فیر ایہ گل ضروری اے کہ اوہ بس اکّو طریقے نال ایس صورت وچ ہو سکدا اے۔ وجہ ویکھن لئی بریکٹاں توں بناں اوہ ٹرم فیر ویکھو جیہڑے اسیں پہلاں تعریف کیتے سن۔ اوہناں لئی ایس ورگی «تعریف» ایہ ہووے گی:
\[
  g(t) =
  \begin{cases}
    0 & \text{ جے $t$ اک اکھر اے}\\
   \max(g(s_1), g(s_2)) + 1 & \text{ جے $t = s_1 \circ s_2$۔}
  \end{cases}
\]
ہُن بریکٹاں توں بناں ٹرم $\mathrm{a} \circ \mathrm{b} \circ
\mathrm{c} \circ \mathrm{d}$ ویکھو۔ ایہنوں کئی طرحاں نال پڑھ سکدے آں، مثلاً $s_1 \circ s_2$ وجوں، جتھے
\begin{align*}
  s_1 & = \mathrm{a} \text{ تے} &
  s_2 & = \mathrm{b} \circ \mathrm{c} \circ \mathrm{d},
\intertext{یا $r_1 \circ r_2$ وجوں، جتھے}
  r_1 & = \mathrm{a} \circ b \text{ تے} &
  r_2 &= \mathrm{c} \circ \mathrm{d}.
\end{align*}
پہلی پڑھت مطابق $g$ دا حساب لائیے تاں ایہ ملے گا:
\begin{align*}
  g(s_1 \circ s_2) & =
  \max(g(\mathrm{a}), g(\mathrm{b} \circ \mathrm{c} \circ \mathrm{d})) + 1 =\\ & =
  \max(0,2) + 1 = 3
  \intertext{جدکہ دوجی پڑھت مطابق سانوں ایہ ملدا اے}
  g(r_1 \circ r_2) & =
  \max(g(\mathrm{a} \circ \mathrm{b}), g(\mathrm{c} \circ \mathrm{d})) + 1=\\ &
  = \max(1,1) + 1 = 2
\end{align*}
پر فنکشن نوں ہر واری اکّو قدر دینی پَندی اے۔ سو~$g$ دی ساڈی «تعریف» کوئی فنکشن تعریف نہیں کردی۔

\begin{prob}
  فنکشن $l$ دی استقرائی تعریف دیو، جتھے $l(t)$، ٹھیک ٹرم~$t$ وچ نشاناں دی گنتی اے۔
\end{prob}

\begin{prob}
  ٹھیک ٹرماں $t$ اُتے ساختی استقرا نال ثابت کرو کہ $f(t) < l(t)$ (جتھے $l(t)$،~$t$ وچ نشاناں دی گنتی اے تے $f(t)$، $t$ دی ڈونگھائی اے جیہدی تعریف \olref[mth][ind][rel]{defn:depth} وچ دتی اے)۔
\end{prob}

\end{document}
